\scanpage{25}
\mspage{12}
qui fait une restriction du groupe structural à $V_J(\mathbb{F}_2)$, d'où $\bigwedge_{i\in J}D_i\simeq\underline{1}$, OK.

Il y aurait lieu aussi d'interpréter les $10$ classes d'isomorphie d'espaces homogènes $\ne(1)$ sous $G\simeq\operatorname{Aut}(\text{Cube})$ (correspondant à des \deleted{\unsure{deux}} sous-groupes d'isotropie $\ne G$, définis à conjugaison près) en termes d'\deleted{\unsure{objets}} éléments de structure associés au cube (les sommets, arêtes, faces, ou lesdits mod antipodisme -- i.e. sur le cube \unsure{gauche} -- ou \deleted{\unsure{repères de paires}}\inserted{repères} d'arêtes -- ce qui en fait déjà sept \inserted{ens.} de cardinaux $8,12,6,\ 4,6,3,\ 24$.
Il y en a un de cardinal $2$, quotient de l'ens des sommets (correspondant à l'ens. $\underline{\omega}_I$), un autre d'ordre $12$ correspondant aux diagonales des faces \inserted{(ou codiagonales, si on préfère)} (les stabilisateurs, \deleted{\unsure{de ces éléments d'ordre}}\inserted{correspondent aux él.} d'ordre $2$ de $G$ qui ne sont pas dans $V$), enfin un troisième qui est formé des paires d'éléments distincts de $J$ (i.e. repères de $J$), $J$ étant interprété comme l'ens des ``directions de faces'' \deleted{\unsure{(=}} (ens des faces/antipodisme) ou des directions d'arêtes. Le compte y est !
\margin{position relative de deux sommets à dist. comb. $1$, ou encore de deux arêtes à dist. comb. \unsure{$2$}, \deleted{ou à dist. \unsure{$2$}}}
\deleted{\unsure{$\mathcal{E}$}}
% archive foliation: 24
